%%%%%%%%%%%%%%%%%%%%%part.tex%%%%%%%%%%%%%%%%%%%%%%%%%%%%%%%%%%
% 
% Part I: The Strategic Vision & Foundations
%
%%%%%%%%%%%%%%%%%%%%%%%% Springer Nature %%%%%%%%%%%%%%%%%%%%%%

\begin{partbacktext}
\part{The Strategic Vision \& Foundations}
\noindent In this first part of the book, we move beyond the technical metrics of accuracy and performance to establish the philosophical and strategic foundations of Trustworthy AI. 

We begin by exploring the \textbf{Accuracy Trap} and the necessary shift toward \textbf{Trust-First engineering}. We then delve into the theory of \textbf{Contextual Integrity}, mapping how information should flow within social contexts, and finally address the \textbf{Socio-Technical} nature of trust, emphasizing the critical role of human-in-the-loop (HITL) design. These three chapters form the bedrock upon which the technical defenses of Part II are built.

\begin{casestudy} 
In early 2025, the SkyLink smart-city project in a major Southeast Asian hub deployed an AI-driven traffic optimization system. While technically successful in reducing congestion by 15 percent, the system inadvertently triggered a massive civil protest. The failure was not one of accuracy, but of \textbf{Contextual Integrity}. By merging transportation data with social media geolocation feeds to predict arrival times, the system began identifying the political affiliations of individuals gathering at public squares. This "Contextual Bleed" from a public utility context into a sensitive political one highlighted the catastrophic risk of treating privacy as a binary state, necessitating the Trust-First paradigms established in Part I. 
\end{casestudy} 
\end{partbacktext}